% !TEX program = xelatex
\documentclass[12pt,a4paper]{ctexbook}

% ==================== 页面 ====================
\usepackage[top=2.5cm,bottom=2.5cm,left=2.5cm,right=2.5cm]{geometry}

% ==================== 字体 ====================
\usepackage{fontspec}
\usepackage{iffont}
\IfFontExistsTF{Consolas}{\setmonofont{Consolas}}{\setmonofont{DejaVu Sans Mono}}

% ==================== 颜色 ====================
\usepackage{xcolor}
\definecolor{codebg}{HTML}{F5F5F5}
\definecolor{codeframe}{HTML}{E0E0E0}
\definecolor{cppkeyword}{HTML}{1A1AFF}
\definecolor{cppcomment}{HTML}{2E8B2E}
\definecolor{cppstring}{HTML}{B22222}
\definecolor{linenumgray}{HTML}{999999}
\definecolor{algheadbg}{HTML}{1A3C5E}

% ==================== listings ====================
\usepackage{listings}

\lstdefinestyle{cppstyle}{
  language=C++,
  basicstyle=\small\ttfamily,
  keywordstyle=\color{cppkeyword}\bfseries,
  commentstyle=\color{cppcomment}\itshape,
  stringstyle=\color{cppstring},
  numberstyle=\tiny\color{linenumgray},
  numbers=left,
  frame=single,
  rulecolor=\color{codeframe},
  framesep=6pt,
  breaklines=true,
  tabsize=4,
  columns=flexible,
  keepspaces=true,
  backgroundcolor=\color{codebg},
  showstringspaces=false,
  morekeywords={constexpr,consteval,constinit,decltype,auto,
    concept,requires,co_await,co_yield,co_return,
    noexcept,override,final,nullptr,static_assert,
    template,typename,namespace,using,mutable,volatile,
    explicit,operator,friend,inline,virtual,
    thread_local,alignas,alignof,if,consteval,constexpr},
  deletekeywords={int,char,bool,float,double,long,short,unsigned,signed,void,wchar_t},
  emph={size_t,int32_t,uint32_t,int64_t,uint64_t,
    string,vector,map,set,unique_ptr,shared_ptr,optional,variant,any,
    span,string_view},
  emphstyle=\color{teal!80!black},
  escapeinside={(*@}{@*)},
}
\lstset{style=cppstyle}

% ==================== tcolorbox ====================
\usepackage[most]{tcolorbox}

\newtcolorbox{guideline}[1][]{
  enhanced,breakable,
  colback=blue!5!white,colframe=blue!75!black,
  fonttitle=\bfseries,title=要点,#1,
  boxed title style={colback=blue!75!black},
  attach boxed title to top left={yshift=-2mm,xshift=2mm},
  arc=2mm,boxrule=0.8mm,
}
\newtcolorbox{compare}[1][]{
  enhanced,breakable,
  colback=green!3!white,colframe=green!50!black,
  fonttitle=\bfseries,title=对比,#1,
  boxed title style={colback=green!50!black},
  attach boxed title to top left={yshift=-2mm,xshift=2mm},
  arc=2mm,boxrule=0.8mm,
}
\newtcolorbox{warning}[1][]{
  enhanced,breakable,
  colback=red!5!white,colframe=red!75!black,
  fonttitle=\bfseries,title=常见陷阱,#1,
  boxed title style={colback=red!75!black},
  attach boxed title to top left={yshift=-2mm,xshift=2mm},
  arc=2mm,boxrule=0.8mm,
}
\newtcolorbox{note}[1][]{
  enhanced,breakable,
  colback=orange!5!white,colframe=orange!80!black,
  fonttitle=\bfseries,title=注意,#1,
  boxed title style={colback=orange!80!black},
  attach boxed title to top left={yshift=-2mm,xshift=2mm},
  arc=2mm,boxrule=0.8mm,
}
\newtcolorbox{bestpractice}[1][]{
  enhanced,breakable,
  colback=purple!5!white,colframe=purple!60!black,
  fonttitle=\bfseries,title=最佳实践,#1,
  boxed title style={colback=purple!60!black},
  attach boxed title to top left={yshift=-2mm,xshift=2mm},
  arc=2mm,boxrule=0.8mm,
}

% ==================== 章节标题 ====================
\usepackage{titlesec}
\usepackage{fancyhdr}
\usepackage{graphicx}
\usepackage{amssymb}
\usepackage{amsmath}
\usepackage{tabularx}
\usepackage{booktabs}
\usepackage{longtable}
\usepackage{array}
\usepackage{multirow}
\usepackage{enumitem}
\usepackage{seqsplit}
\usepackage{float}

\titleformat{\chapter}[display]
  {\normalfont\huge\bfseries\color{algheadbg}}
  {\chaptertitlename\ \thechapter}{20pt}{\Huge}
\titleformat{\section}
  {\normalfont\Large\bfseries\color{blue!70!black}}
  {\thesection}{1em}{}
\titleformat{\subsection}
  {\normalfont\large\bfseries\color{blue!60!black}}
  {\thesubsection}{1em}{}

\pagestyle{fancy}
\fancyhf{}
\fancyhead[LE,RO]{\thepage}
\fancyhead[RE]{\leftmark}
\fancyhead[LO]{\rightmark}

\setlist[itemize]{leftmargin=2em,itemsep=2pt}
\setlist[enumerate]{leftmargin=2em,itemsep=2pt}

% ==================== 超链接 ====================
\usepackage{hyperref}
\hypersetup{
  colorlinks=true,
  linkcolor=blue!70!black,
  urlcolor=blue!60!black,
  bookmarks=true,
  pdfauthor={C++ Reference},
  pdftitle={C++ new/delete 运算符重载},
}

% ==================== 自定义命令 ====================
\newcommand{\cpp}[1]{C++#1}
\newcommand{\Fn}[1]{\texttt{#1()}}
\newcommand{\Tp}[1]{\texttt{#1}}

% ====================================================================
\begin{document}

\frontmatter

\begin{titlepage}
\centering
\vspace*{4cm}
{\Huge\bfseries C++ new/delete 运算符重载\par}
\vspace{1cm}
{\LARGE 内存调试、裸机、外设分配与自定义分配器\par}
\vfill
{\large \today\par}
\end{titlepage}

\tableofcontents

\mainmatter

% ================================================================
\chapter{operator new/delete 家族概览}
% ================================================================

C++ 将内存分配拆分为\textbf{原始内存获取}与\textbf{对象构造}两个层次。\texttt{operator new}/\texttt{operator delete} 仅负责获取与释放原始内存，\texttt{new}/\texttt{delete} 表达式则自动调用构造与析构函数。理解这一区分是掌握全部重载机制的前提。

\section{六个重载点}

标准库在 \texttt{<new>} 头文件中声明了六对全局重载函数。每一对都包含
\texttt{nothrow} 版本和可能抛出 \texttt{std::bad\_alloc} 的版本。

\begin{lstlisting}
// === 单对象分配 ===
void* operator new  (std::size_t count);
void* operator new  (std::size_t count,
                     const std::nothrow_t&) noexcept;

// === 单对象释放 ===
void  operator delete(void* ptr) noexcept;
void  operator delete(void* ptr,
                      std::size_t count) noexcept;

// === 数组分配 ===
void* operator new[](std::size_t count);
void* operator new[](std::size_t count,
                     const std::nothrow_t&) noexcept;

// === 数组释放 ===
void  operator delete[](void* ptr) noexcept;
void  operator delete[](void* ptr,
                        std::size_t count) noexcept;
\end{lstlisting}

\begin{note}
\cpp{14} 起，\texttt{operator delete} 增加了带 \texttt{size\_t} 参数的
sized deallocation 版本。编译器在知道对象大小时会自动调用带尺寸的版本，
允许分配器进行更高效的回收。
\end{note}

\section{new 表达式与 operator new 的关系}

以下伪代码展示了编译器对 \texttt{new} 表达式的展开过程：

\begin{lstlisting}
// 程序员写的代码:
Widget* w = new Widget(arg1, arg2);

// 编译器展开为:
void* raw = operator new(sizeof(Widget));   // 1. 分配原始内存
try {
    w = new(raw) Widget(arg1, arg2);        // 2. placement new 构造
} catch (...) {
    operator delete(raw);                   // 3. 构造失败则释放
    throw;
}
\end{lstlisting}

\texttt{delete} 表达式的展开是镜像过程：先调用析构函数，再调用
\texttt{operator delete} 释放内存。

\section{数组版本的额外开销}

\texttt{new[]}/ \texttt{delete[]} 的实现通常会在分配的内存块头部写入
数组元素数量（cookie），以便 \texttt{delete[]} 知道需要调用多少次析构函数。
这意味着：

\begin{lstlisting}
Widget* arr = new Widget[10];
// 实际分配: sizeof(size_t) + sizeof(Widget) * 10
// cookie 存放元素数量 10
// arr 指向 cookie 之后第一个 Widget 的地址

delete[] arr;  // 读 cookie, 调用 10 次析构, 释放整块
\end{lstlisting}

\begin{warning}
对 \texttt{new[]} 的结果使用 \texttt{delete}（不带方括号）会导致
未定义行为：cookie 不会被读取，析构函数调用次数错误，且释放的指针
可能不指向分配时的实际起始地址。
\end{warning}

\section{Placement new}

Placement new 是 \texttt{operator new} 的一个特殊重载，接受额外的
地址参数，在已有内存上构造对象。它不分配任何内存：

\begin{lstlisting}
#include <new>

alignas(Widget) char buffer[sizeof(Widget)];
Widget* w = new (buffer) Widget(42);    // 在 buffer 上构造

w->~Widget();                           // 手动析构
// 不调用 operator delete, 因为内存不是 new 分配的
\end{lstlisting}

Placement new 在嵌入式系统、内存池和自定义分配器中广泛使用。
对应的析构必须手动调用析构函数，不能使用 \texttt{delete} 表达式。

% ================================================================
\chapter{全局重载：内存调试与泄漏检测}
% ================================================================

全局重载 \texttt{operator new} 和 \texttt{operator delete} 可以拦截
程序中所有的动态内存分配。这是实现内存调试、泄漏检测和分配统计的基础手段。

\section{基本原理}

在翻译单元中定义与全局版本相同签名的函数即可替换默认实现。
链接器会使用自定义版本覆盖标准库的默认版本：

\begin{lstlisting}
#include <cstdlib>
#include <cstdio>
#include <new>

void* operator new(std::size_t count) {
    void* ptr = std::malloc(count);
    if (!ptr) throw std::bad_alloc();
    std::printf("[ALLOC] %zu bytes at %p\n", count, ptr);
    return ptr;
}

void operator delete(void* ptr) noexcept {
    std::printf("[FREE]  %p\n", ptr);
    std::free(ptr);
}

// 必须同时重载 nothrow 版本和数组版本
void* operator new(std::size_t count,
                   const std::nothrow_t&) noexcept {
    try { return operator new(count); }
    catch (...) { return nullptr; }
}

void operator delete(void* ptr, std::size_t) noexcept {
    operator delete(ptr);
}

void* operator new[](std::size_t count) {
    return operator new(count);
}

void operator delete[](void* ptr) noexcept {
    operator delete(ptr);
}

void* operator new[](std::size_t count,
                     const std::nothrow_t&) noexcept {
    return operator new(count, std::nothrow);
}

void operator delete[](void* ptr, std::size_t) noexcept {
    operator delete(ptr);
}
\end{lstlisting}

\begin{warning}
全局重载时，自定义 \texttt{operator new} 不能内部使用 \texttt{new} 表达式，
否则会导致无限递归。必须使用 \texttt{malloc}/\texttt{free} 或
\texttt{mmap}/\texttt{munmap} 等底层 API 进行实际内存分配。
\end{warning}

\section{内存分配追踪}

通过在全局重载中记录每次分配的信息（地址、大小），可以构建一个
完整的内存使用图景。以下实现追踪所有活跃分配：

\begin{lstlisting}
#include <unordered_map>
#include <mutex>
#include <cstdio>
#include <cstdlib>
#include <new>

namespace {
    struct AllocInfo { std::size_t size; };
    std::unordered_map<void*, AllocInfo> g_allocs;
    std::mutex g_mutex;
    std::size_t g_total = 0;
    std::size_t g_peak  = 0;
    std::size_t g_count = 0;
}

void* operator new(std::size_t count) {
    void* ptr = std::malloc(count);
    if (!ptr) throw std::bad_alloc();
    {
        std::lock_guard lock(g_mutex);
        g_allocs[ptr] = {count};
        g_total += count;
        g_count++;
        if (g_total > g_peak) g_peak = g_total;
    }
    return ptr;
}

void operator delete(void* ptr) noexcept {
    if (!ptr) return;
    {
        std::lock_guard lock(g_mutex);
        auto it = g_allocs.find(ptr);
        if (it != g_allocs.end()) {
            g_total -= it->second.size;
            g_allocs.erase(it);
        }
    }
    std::free(ptr);
}
\end{lstlisting}

程序退出时，\texttt{g\_allocs} 中残留的条目即为泄漏的内存。
可以注册 \texttt{atexit} 回调在程序结束时打印泄漏报告：

\begin{lstlisting}
void report_leaks() {
    std::lock_guard lock(g_mutex);
    if (g_allocs.empty()) {
        std::printf("[MEM] No leaks detected.\n");
        return;
    }
    std::printf("[MEM] %zu leaks, %zu bytes total\n",
                g_allocs.size(), g_total);
    for (auto& [ptr, info] : g_allocs) {
        std::printf("  %p : %zu bytes\n", ptr, info.size);
    }
    std::printf("[MEM] Peak usage: %zu bytes\n", g_peak);
}

// 全局对象构造函数中注册
static struct LeakReporter {
    LeakReporter() { std::atexit(report_leaks); }
} g_reporter;
\end{lstlisting}

\section{源码位置标记}

结合 \cpp{20} 的 \texttt{std::source\_location}，可以在分配时记录
调用位置。定义带额外参数的 placement new：

\begin{lstlisting}
#include <source_location>

void* operator new(std::size_t count,
                   const std::source_location& loc) {
    void* ptr = std::malloc(count);
    if (!ptr) throw std::bad_alloc();
    std::printf("[ALLOC] %zu bytes at %p from %s:%d\n",
                count, ptr, loc.file_name(), loc.line());
    return ptr;
}

// 宏: 替换 new 表达式为携带位置信息的版本
#define TRACKED_NEW new(std::source_location::current())

// 使用:
// Widget* w = TRACKED_NEW Widget(42);
// 输出: [ALLOC] 24 bytes at 0x... from main.cpp:42
\end{lstlisting}

\begin{note}[title=Placement delete]
带额外参数的 placement \texttt{operator new} 不会替换全局版本，
而是作为重载解析的额外候选。\texttt{delete} 表达式仍然调用标准的全局
\texttt{operator delete}，因此泄漏检测逻辑只需在分配端记录即可。
\end{note}

\section{裸机环境：无 malloc 的系统}

在裸机（bare-metal）嵌入式系统中，标准库的 \texttt{malloc}/\texttt{free}
可能根本不存在。此时必须全局重载 \texttt{operator new}，从预定义的内存池
中分配内存。

\subsection{静态内存池实现}

\begin{lstlisting}
#include <cstddef>
#include <new>

// 预分配的静态内存池（位于特定 RAM 段）
static constexpr std::size_t POOL_SIZE = 64 * 1024;
alignas(16) static char g_pool[POOL_SIZE];
static std::size_t g_offset = 0;

void* operator new(std::size_t count) {
    // 对齐到 16 字节边界
    std::size_t aligned = (g_offset + 15) & ~std::size_t(15);
    if (aligned + count > POOL_SIZE)
        throw std::bad_alloc();
    void* ptr = &g_pool[aligned];
    g_offset = aligned + count;
    return ptr;
}

void* operator new(std::size_t count,
                   const std::nothrow_t&) noexcept {
    try { return operator new(count); }
    catch (...) { return nullptr; }
}

// 裸机线性分配器: 不支持释放
void operator delete(void*) noexcept { /* no-op */ }
void operator delete(void*, std::size_t) noexcept {}
void* operator new[](std::size_t count) {
    return operator new(count);
}
void operator delete[](void*) noexcept { /* no-op */ }
void operator delete[](void*, std::size_t) noexcept {}
\end{lstlisting}

\subsection{链接脚本配合}

在实际嵌入式项目中，内存池通常由链接脚本定义在特定的内存段中。
C++ 代码通过 \texttt{extern} 符号引用链接器提供的地址：

\begin{lstlisting}
extern "C" {
    extern char __heap_start;
    extern char __heap_end;
}

void* operator new(std::size_t count) {
    static char* heap_ptr = &__heap_start;
    auto aligned = reinterpret_cast<char*>(
        (reinterpret_cast<std::uintptr_t>(heap_ptr) + 15)
        & ~std::uintptr_t(15));
    if (aligned + count > &__heap_end)
        throw std::bad_alloc();
    heap_ptr = aligned + count;
    return aligned;
}
\end{lstlisting}

\begin{guideline}
裸机环境的全局 \texttt{operator new} 实现要点：
\begin{itemize}
    \item 使用静态数组或链接脚本定义的内存段，避免依赖运行时库。
    \item 对齐通常为 8 或 16 字节，取决于目标平台的指针大小和 SIMD 需求。
    \item 线性分配器（bump allocator）最简单，但不支持释放。
    \item 如需释放，需实现简单的空闲链表或 buddy 分配器。
    \item 在多线程裸机环境中，\texttt{operator new} 必须加锁。
\end{itemize}
\end{guideline}

\section{外设内存分配}

某些硬件架构要求特定数据必须位于特定内存区域。例如，DMA 缓冲区需要在
物理连续的内存中，GPU 数据需要在显存中，共享内存需要在多核可见的区域。
全局重载 \texttt{operator new} 可以实现对外设内存的透明分配。

\subsection{DMA 缓冲区分配}

\begin{lstlisting}
#include <cstddef>
#include <cstdint>
#include <new>

// 硬件 API: 分配物理连续的 DMA 内存
extern "C" void* dma_alloc(std::size_t size,
                            std::uintptr_t* phys_addr);
extern "C" void  dma_free(void* virt_addr, std::size_t size);

// 标记结构
struct DmaTag {};
inline constexpr DmaTag dma_tag{};

// placement new: 分配 DMA 可用内存
void* operator new(std::size_t count, DmaTag) {
    std::uintptr_t phys = 0;
    void* ptr = dma_alloc(count, &phys);
    if (!ptr) throw std::bad_alloc();
    return ptr;
}

// placement delete: 构造失败时调用
void operator delete(void* ptr, DmaTag) noexcept {
    dma_free(ptr, 0);
}

// 使用: auto* buf = new (dma_tag) DmaBuffer(4096);
\end{lstlisting}

\begin{note}[title=Placement delete 的触发时机]
placement \texttt{operator delete} 只在对应的 placement \texttt{new}
表达式中对象构造抛出异常时才会被调用。如果构造成功，程序员需要手动
调用析构函数并自行释放内存。
\end{note}

% ================================================================
\chapter{类级别重载}
% ================================================================

类级别的 \texttt{operator new} 和 \texttt{operator delete} 作为类的静态
成员函数，仅影响该类对象的内存分配。它不影响标准库容器内部的元素分配
（那由 allocator 控制），只在使用 \texttt{new} 表达式创建该类的对象时生效。

\section{固定大小的对象池}

对于频繁创建和销毁的同类型对象，类级别的 operator new 可以实现
对象池，避免反复调用系统分配器：

\begin{lstlisting}
#include <cstddef>
#include <new>
#include <array>

class Particle {
public:
    float x, y, z;
    float vx, vy, vz;
    int   lifetime;

    static void* operator new(std::size_t) {
        if (free_list_) {
            Node* n = free_list_;
            free_list_ = n->next;
            return n;
        }
        return ::operator new(sizeof(Particle));
    }

    static void operator delete(void* ptr) noexcept {
        Node* n = static_cast<Node*>(ptr);
        n->next = free_list_;
        free_list_ = n;
    }

private:
    struct Node { Node* next; };
    static inline Node* free_list_ = nullptr;
    static inline std::array<Particle, 1024> pool_{};

    struct PoolInit {
        PoolInit() {
            for (auto& p : pool_) {
                auto* n = reinterpret_cast<Node*>(&p);
                n->next = free_list_;
                free_list_ = n;
            }
        }
    };
    static inline PoolInit init_{};
};
\end{lstlisting}

使用方式与普通 \texttt{new} 完全相同，但内存来自预分配的对象池，
分配和释放均为 $O(1)$ 时间复杂度：

\begin{lstlisting}
Particle* p = new Particle{
    1.0f, 0.0f, 0.0f, 0.5f, 0.0f, 0.0f, 100};
delete p;   // 归还到对象池
\end{lstlisting}

\section{GPU 向量：分配到显存的完整示例}

以下实现一个 \texttt{GpuVector} 类，其数据存储在 GPU 可访问的内存中。
类级别的 \texttt{operator new} 将对象本身（包括指针和元数据）分配到
主机内存，而实际的数据缓冲区通过 GPU API 分配到显存。

\subsection{GPU 内存管理接口}

首先定义与 GPU 驱动交互的抽象接口：

\begin{lstlisting}
#include <cstddef>
#include <cstdint>
#include <new>

namespace gpu {

enum class MemoryType {
    DeviceLocal,       // 显存, GPU 高速访问
    HostVisible,       // 主机可见, CPU 可映射
    UnifiedMemory      // 统一内存, CPU/GPU 共享
};

void* allocate(std::size_t size, MemoryType type);
void  deallocate(void* ptr);
void  upload(void* dst, const void* src, std::size_t size);
void  download(void* dst, const void* src, std::size_t size);

} // namespace gpu
\end{lstlisting}

\subsection{GpuVector 类定义}

\begin{lstlisting}
#include <cstddef>
#include <new>
#include <initializer_list>

template <typename T>
class GpuVector {
public:
    GpuVector() = default;

    explicit GpuVector(std::size_t count) {
        reserve(count);
        size_ = count;
    }

    GpuVector(std::initializer_list<T> init) {
        reserve(init.size());
        size_ = init.size();
        gpu::upload(gpu_data_, init.begin(),
                    size_ * sizeof(T));
    }

    ~GpuVector() {
        if (gpu_data_) gpu::deallocate(gpu_data_);
        delete[] staging_;
    }

    GpuVector(const GpuVector&) = delete;
    GpuVector& operator=(const GpuVector&) = delete;

    GpuVector(GpuVector&& o) noexcept
        : gpu_data_(o.gpu_data_),
          staging_(o.staging_),
          size_(o.size_),
          capacity_(o.capacity_) {
        o.gpu_data_ = nullptr;
        o.staging_  = nullptr;
        o.size_     = 0;
        o.capacity_ = 0;
    }

    GpuVector& operator=(GpuVector&& o) noexcept {
        if (this != &o) {
            if (gpu_data_) gpu::deallocate(gpu_data_);
            delete[] staging_;
            gpu_data_  = o.gpu_data_;
            staging_   = o.staging_;
            size_      = o.size_;
            capacity_  = o.capacity_;
            o.gpu_data_ = nullptr;
            o.staging_  = nullptr;
            o.size_     = 0;
            o.capacity_ = 0;
        }
        return *this;
    }

    T read(std::size_t index) const {
        sync_to_host();
        return staging_[index];
    }

    void write(std::size_t index, const T& value) {
        ensure_staging();
        staging_[index] = value;
        gpu::upload(
            static_cast<char*>(gpu_data_)
                + index * sizeof(T),
            &value, sizeof(T));
    }

    void push_back(const T& value) {
        if (size_ >= capacity_) {
            std::size_t new_cap = capacity_ == 0
                ? 8 : capacity_ * 2;
            reserve(new_cap);
        }
        write(size_++, value);
    }

    std::size_t size() const { return size_; }
    bool empty() const { return size_ == 0; }

    // GPU 指针, 供 shader / kernel 直接使用
    T* gpu_data() {
        return static_cast<T*>(gpu_data_);
    }
    const T* gpu_data() const {
        return static_cast<const T*>(gpu_data_);
    }

private:
    void* gpu_data_  = nullptr;
    T*    staging_   = nullptr;
    std::size_t size_     = 0;
    std::size_t capacity_ = 0;

    void reserve(std::size_t new_cap) {
        if (new_cap <= capacity_) return;
        void* new_gpu = gpu::allocate(
            new_cap * sizeof(T),
            gpu::MemoryType::DeviceLocal);
        if (gpu_data_ && size_ > 0) {
            gpu::upload(new_gpu, gpu_data_,
                        size_ * sizeof(T));
            gpu::deallocate(gpu_data_);
        }
        delete[] staging_;
        staging_   = new T[new_cap];
        gpu_data_  = new_gpu;
        capacity_  = new_cap;
        if (size_ > 0) {
            gpu::download(staging_, gpu_data_,
                          size_ * sizeof(T));
        }
    }

    void ensure_staging() {
        if (!staging_ && capacity_ > 0) {
            staging_ = new T[capacity_];
            if (gpu_data_ && size_ > 0) {
                gpu::download(staging_, gpu_data_,
                              size_ * sizeof(T));
            }
        }
    }

    void sync_to_host() const {
        if (gpu_data_ && staging_ && size_ > 0) {
            gpu::download(staging_, gpu_data_,
                          size_ * sizeof(T));
        }
    }
};
\end{lstlisting}

\subsection{类级别 operator new 重载}

\texttt{GpuVector} 类级别的 \texttt{operator new} 控制的是
\texttt{GpuVector} 对象本身（指针 + 元数据）的分配位置，而非其管理的
数据缓冲区。如果需要将对象本身也放到特殊内存中（例如嵌入式系统的
特定 SRAM 段），可以重载类级别的 operator new：

\begin{lstlisting}
template <typename T>
class GpuVector {
public:
    // ... (同上) ...

    // 类级别重载: 将对象本身分配到 GPU 统一内存
    static void* operator new(std::size_t count) {
        void* ptr = gpu::allocate(
            count, gpu::MemoryType::UnifiedMemory);
        if (!ptr) throw std::bad_alloc();
        return ptr;
    }

    static void operator delete(void* ptr) noexcept {
        gpu::deallocate(ptr);
    }

    static void* operator new[](std::size_t count) {
        return operator new(count);
    }

    static void operator delete[](void* ptr) noexcept {
        operator delete(ptr);
    }
};
\end{lstlisting}

此时 \texttt{new GpuVector<float>(100)} 会将 \texttt{GpuVector}
对象本身分配到统一内存，而对象内部的 \texttt{gpu\_data\_} 指针指向
设备本地显存。两层分配各司其职。

\subsection{使用示例}

\begin{lstlisting}
int main() {
    // 对象本身在统一内存, 数据在显存
    auto* vec = new GpuVector<float>(1024);

    // 写入数据（自动上传到 GPU）
    for (std::size_t i = 0; i < 1024; i++) {
        vec->write(i, static_cast<float>(i) * 0.1f);
    }

    // 将 GPU 指针传给计算 shader
    // kernel_launch(vec->gpu_data(), vec->size());

    // 读取回主机
    float val = vec->read(42);

    delete vec;
    return 0;
}
\end{lstlisting}

\begin{guideline}
\texttt{GpuVector} 的设计要点：
\begin{itemize}
    \item 数据缓冲区通过 GPU API 分配到显存，这是构造函数内部的逻辑，
          与 \texttt{operator new} 无关。
    \item 类级别 \texttt{operator new} 控制的是 \texttt{GpuVector}
          对象本身（约 32 字节的指针和整数）的分配位置。
    \item 如果需要 CPU 端频繁读取，应维护主机端的 staging buffer
          作为缓存，减少 GPU--CPU 数据传输。
    \item \texttt{gpu\_data()} 返回裸指针供 GPU shader 直接使用，
          避免任何 CPU 端的间接访问。
\end{itemize}
\end{guideline}

% ================================================================
\chapter{与自定义分配器的区别}
% ================================================================

类级别的 \texttt{operator new} 重载和 STL 自定义分配器
（\texttt{Allocator}）都涉及"控制内存从哪里来"，但它们作用的层次、
生效范围和适用场景完全不同。

\section{作用层次对比}

\texttt{operator new} 工作在\textbf{表达式}层次：每当使用 \texttt{new}
表达式创建对象时，编译器调用该类的 \texttt{operator new} 获取内存。
而 \texttt{Allocator} 工作在\textbf{容器内部}：标准库容器使用分配器
来分配其内部元素存储，与对象本身的分配方式无关。

\begin{compare}[title=operator new vs Allocator]
\textbf{operator new 重载}：
\begin{itemize}
    \item 作用范围：\texttt{new T} 和 \texttt{delete p} 表达式
    \item 控制目标：对象 \texttt{T} 本身的内存来源
    \item 触发条件：使用 \texttt{new}/\texttt{delete} 创建/销毁对象
    \item 典型场景：对象池、嵌入式裸机、GPU 统一内存
\end{itemize}

\textbf{自定义 Allocator}：
\begin{itemize}
    \item 作用范围：标准库容器内部的元素存储
    \item 控制目标：\texttt{vector<T>} 中 \texttt{T} 元素的内存来源
    \item 触发条件：容器扩容、元素插入/删除
    \item 典型场景：内存池分配器、共享内存分配器、GPU 数据分配器
\end{itemize}
\end{compare}

\section{具体差异演示}

以下代码展示了两种机制各自影响的位置：

\begin{lstlisting}
// 场景 1: operator new 重载
class Widget {
public:
    static void* operator new(std::size_t);
    static void  operator delete(void*) noexcept;
    int data[16];
};

Widget* w = new Widget;           // Widget::operator new
auto* v = new std::vector<Widget>;
v->push_back(Widget{});
// push_back 内部分配 Widget 时:
// 使用 std::allocator<Widget>::allocate
// 不调用 Widget::operator new
\end{lstlisting}

\begin{lstlisting}
// 场景 2: 自定义 Allocator
template <typename T>
struct GpuAllocator {
    using value_type = T;
    T* allocate(std::size_t n) {
        return static_cast<T*>(gpu::allocate(
            n * sizeof(T),
            gpu::MemoryType::DeviceLocal));
    }
    void deallocate(T* p, std::size_t) noexcept {
        gpu::deallocate(p);
    }
};

// 容器内部的元素存储在 GPU 显存中
std::vector<float, GpuAllocator<float>> data;
data.push_back(3.14f);    // GpuAllocator::allocate

// 但 vector 对象本身仍在堆上
auto* pv = new std::vector<float, GpuAllocator<float>>;
// ^ 使用全局 ::operator new 分配 vector 控制块
\end{lstlisting}

\section{协同使用}

在实际的 GPU 计算场景中，两种机制可以协同工作：
\texttt{operator new} 将对象分配到 GPU 可见的统一内存，
自定义分配器将容器元素分配到设备本地显存。

\begin{lstlisting}
// 自定义分配器: 元素分配到 GPU 显存
template <typename T>
struct DeviceAllocator {
    using value_type = T;
    T* allocate(std::size_t n) {
        return static_cast<T*>(gpu::allocate(
            n * sizeof(T),
            gpu::MemoryType::DeviceLocal));
    }
    void deallocate(T* p, std::size_t) noexcept {
        gpu::deallocate(p);
    }
    bool operator==(const DeviceAllocator&) const {
        return true;
    }
    bool operator!=(const DeviceAllocator&) const {
        return false;
    }
};

struct ParticleSystem {
    // 粒子数据在显存
    std::vector<Particle, DeviceAllocator<Particle>>
        particles;

    // 对象本身在 GPU 可见的统一内存
    static void* operator new(std::size_t sz) {
        return gpu::allocate(sz,
            gpu::MemoryType::UnifiedMemory);
    }
    static void operator delete(void* ptr) noexcept {
        gpu::deallocate(ptr);
    }
};

auto* sys = new ParticleSystem;          // 统一内存
sys->particles.resize(100000);           // 显存
// GPU 内核直接读取 sys->particles.data()
\end{lstlisting}

\section{PMR 多态分配器}

\cpp{17} 引入的 \texttt{std::pmr::polymorphic\_allocator} 将分配策略
从编译时模板参数变为运行时选择。它基于 \texttt{memory\_resource} 接口，
允许在不改变容器类型的情况下切换底层分配策略：

\begin{lstlisting}
#include <memory_resource>

class GpuMemoryResource
    : public std::pmr::memory_resource {
protected:
    void* do_allocate(std::size_t bytes,
                      std::size_t /*alignment*/) override {
        return gpu::allocate(bytes,
            gpu::MemoryType::DeviceLocal);
    }
    void do_deallocate(void* p, std::size_t,
                       std::size_t) override {
        gpu::deallocate(p);
    }
    bool do_is_equal(
        const memory_resource& o) const noexcept override {
        return this == &o;
    }
};

// 使用
GpuMemoryResource gpu_res;
std::pmr::vector<float> data(&gpu_res);
// data 的内部存储从 gpu_res 分配（显存）

// 切换到主机内存, 类型不变
std::pmr::vector<float> host_data(
    std::pmr::new_delete_resource());
\end{lstlisting}

\begin{bestpractice}[title=选择正确的机制]
根据需求选择合适的内存控制方式：
\begin{itemize}
    \item \textbf{全局 operator new}：需要拦截所有分配（内存调试、泄漏检测），
          或在裸机环境中完全替代 \texttt{malloc}。
    \item \textbf{类级别 operator new}：特定类型的对象需要特殊分配位置
          （对象池、统一内存），且主要通过 \texttt{new}/\texttt{delete} 创建。
    \item \textbf{自定义 Allocator}：标准库容器的内部元素需要特殊分配
          （显存、共享内存、NUMA 内存），容器类型在编译时确定。
    \item \textbf{PMR 分配器}：需要在运行时动态切换分配策略，
          或需要容器在不同分配器之间互操作。
\end{itemize}
\end{bestpractice}

\section{完整决策流程}

面对一个新的内存管理需求时，按以下逻辑决策：

\begin{enumerate}
    \item \textbf{需求是否涉及所有类型？}——如果是（如全局内存调试），
          使用全局 \texttt{operator new} 重载。
    \item \textbf{需求仅针对特定类型的对象？}——如果是（如对象池），
          使用类级别 \texttt{operator new}。
    \item \textbf{需求是否针对容器内部的元素存储？}——如果是
          （如 GPU 显存上的 \texttt{vector}），使用自定义 Allocator。
    \item \textbf{分配策略是否需要运行时切换？}——如果是，
          使用 PMR \texttt{memory\_resource}。
    \item \textbf{需要同时控制对象和元素？}——组合使用
          类级别 \texttt{operator new} + 自定义 Allocator。
\end{enumerate}


% ================================================================
\chapter{CUDA 实战：两种封装方式对比}
% ================================================================

本章以 CUDA 为具体平台，演示如何封装一个 GPU 向量。
同一需求可以用两种截然不同的技术路线实现：
重载 \texttt{operator new/delete}，或者编写自定义 \texttt{Allocator}。
对比两者有助于加深理解上一章的理论区分。

\section{CUDA 内存管理基础}

CUDA 运行时提供以下核心 API：

\begin{lstlisting}
#include <cuda_runtime.h>

// 分配/释放设备内存（显存）
cudaError_t cudaMalloc(void** devPtr, size_t size);
cudaError_t cudaFree(void* devPtr);

// 主机与设备之间的数据传输
cudaError_t cudaMemcpy(void* dst, const void* src,
                       size_t count,
                       cudaMemcpyKind kind);
// kind: cudaMemcpyHostToDevice / cudaMemcpyDeviceToHost

// 分配统一内存（CPU/GPU 均可访问）
cudaError_t cudaMallocManaged(void** devPtr,
                              size_t size);
\end{lstlisting}

以下封装中，\texttt{CudaVector} 的数据缓冲区始终通过
\texttt{cudaMalloc} 分配到显存，而主机端维护一个 staging 副本
用于 CPU 侧的元素访问。

\section{方式一：类级别 operator new/delete}

这种方式通过重载 \texttt{CudaVector} 类自身的分配函数，
将对象控制块（指针 + 尺寸元数据）也放到 CUDA 统一内存中，
使得 CPU 和 GPU 代码都可以直接访问 \texttt{CudaVector} 的成员变量。

\subsection{完整实现}

\begin{lstlisting}
#include <cuda_runtime.h>
#include <cstddef>
#include <new>
#include <stdexcept>
#include <initializer_list>

template <typename T>
class CudaVector {
public:
    CudaVector() = default;

    explicit CudaVector(std::size_t n) {
        resize(n);
    }

    CudaVector(std::size_t n, const T& val) {
        resize(n);
        for (std::size_t i = 0; i < n; i++)
            staging_[i] = val;
        upload_all();
    }

    CudaVector(std::initializer_list<T> init) {
        resize(init.size());
        std::size_t i = 0;
        for (auto& v : init) staging_[i++] = v;
        upload_all();
    }

    ~CudaVector() {
        if (device_ptr_) cudaFree(device_ptr_);
        delete[] staging_;
    }

    // 禁止拷贝
    CudaVector(const CudaVector&) = delete;
    CudaVector& operator=(const CudaVector&) = delete;

    // 允许移动
    CudaVector(CudaVector&& o) noexcept
        : device_ptr_(o.device_ptr_),
          staging_(o.staging_),
          size_(o.size_),
          capacity_(o.capacity_) {
        o.device_ptr_ = nullptr;
        o.staging_    = nullptr;
        o.size_       = 0;
        o.capacity_   = 0;
    }

    CudaVector& operator=(CudaVector&& o) noexcept {
        if (this != &o) {
            if (device_ptr_) cudaFree(device_ptr_);
            delete[] staging_;
            device_ptr_ = o.device_ptr_;
            staging_    = o.staging_;
            size_       = o.size_;
            capacity_   = o.capacity_;
            o.device_ptr_ = nullptr;
            o.staging_    = nullptr;
            o.size_       = 0;
            o.capacity_   = 0;
        }
        return *this;
    }

    // --- 类级别 operator new/delete ---
    // 将 CudaVector 对象本身分配到统一内存
    static void* operator new(std::size_t sz) {
        void* ptr = nullptr;
        auto err = cudaMallocManaged(&ptr, sz);
        if (err != cudaSuccess || !ptr)
            throw std::bad_alloc();
        return ptr;
    }

    static void operator delete(void* ptr) noexcept {
        if (ptr) cudaFree(ptr);
    }

    static void* operator new[](std::size_t sz) {
        return operator new(sz);
    }

    static void operator delete[](void* p) noexcept {
        operator delete(p);
    }

    // --- 元素访问 ---
    T read(std::size_t i) const {
        sync_to_host();
        return staging_[i];
    }

    void write(std::size_t i, const T& val) {
        staging_[i] = val;
        cudaMemcpy(
            device_ptr_ + i, &val, sizeof(T),
            cudaMemcpyHostToDevice);
    }

    void push_back(const T& val) {
        if (size_ >= capacity_)
            grow(capacity_ == 0 ? 8 : capacity_ * 2);
        staging_[size_] = val;
        cudaMemcpy(
            device_ptr_ + size_, &val, sizeof(T),
            cudaMemcpyHostToDevice);
        size_++;
    }

    std::size_t size() const { return size_; }
    bool empty() const { return size_ == 0; }

    // 供 CUDA kernel 直接使用
    T* device_data() { return device_ptr_; }
    const T* device_data() const { return device_ptr_; }

private:
    T*    device_ptr_ = nullptr;   // 显存
    T*    staging_    = nullptr;   // 主机暂存
    std::size_t size_     = 0;
    std::size_t capacity_ = 0;

    void resize(std::size_t n) {
        if (n > capacity_) grow(n);
        size_ = n;
    }

    void grow(std::size_t new_cap) {
        T* new_dev = nullptr;
        auto err = cudaMalloc(&new_dev,
                              new_cap * sizeof(T));
        if (err != cudaSuccess)
            throw std::bad_alloc();

        if (device_ptr_ && size_ > 0) {
            // 设备到设备: 通过主机中转
            cudaMemcpy(new_dev, device_ptr_,
                       size_ * sizeof(T),
                       cudaMemcpyDeviceToDevice);
            cudaFree(device_ptr_);
        }

        T* new_staging = new T[new_cap];
        if (staging_ && size_ > 0) {
            for (std::size_t i = 0; i < size_; i++)
                new_staging[i] = staging_[i];
        }
        delete[] staging_;

        device_ptr_ = new_dev;
        staging_    = new_staging;
        capacity_   = new_cap;
    }

    void upload_all() {
        if (device_ptr_ && staging_ && size_ > 0) {
            cudaMemcpy(device_ptr_, staging_,
                       size_ * sizeof(T),
                       cudaMemcpyHostToDevice);
        }
    }

    void sync_to_host() const {
        if (device_ptr_ && staging_ && size_ > 0) {
            cudaMemcpy(staging_, device_ptr_,
                       size_ * sizeof(T),
                       cudaMemcpyDeviceToHost);
        }
    }
};
\end{lstlisting}

\subsection{使用方式与 CUDA kernel 调用}

\begin{lstlisting}
// CUDA kernel: 每个线程乘 2
__global__ void scale_kernel(float* data, int n) {
    int idx = blockIdx.x * blockDim.x + threadIdx.x;
    if (idx < n) data[idx] *= 2.0f;
}

int main() {
    // new CudaVector: 对象本身在统一内存
    auto* vec = new CudaVector<float>(1024, 1.0f);

    // 启动 kernel，直接操作 vec->device_data()
    int threads = 256;
    int blocks = (1024 + threads - 1) / threads;
    scale_kernel<<<blocks, threads>>>(
        vec->device_data(),
        static_cast<int>(vec->size()));
    cudaDeviceSynchronize();

    // CPU 端读取（自动从显存同步）
    float val = vec->read(0);   // 2.0

    delete vec;   // cudaFree 对象 + 数据
    return 0;
}
\end{lstlisting}

\begin{note}[title=统一内存的代价]
\texttt{cudaMallocManaged} 分配的统一内存由 CUDA 驱动自动在 CPU 和 GPU
之间迁移页面。这简化了编程，但在频繁跨设备访问时会产生 page fault 开销。
对性能敏感的场景，应将对象控制块放在主机内存，仅数据放在显存。
\end{note}

\section{方式二：自定义 CUDA Allocator}

第二种方式保持 \texttt{CudaVector} 作为普通类（对象在主机堆上），
但为 \texttt{std::vector} 编写一个 CUDA 分配器，使容器内部的
元素存储位于显存。

\subsection{CudaAllocator 实现}

\begin{lstlisting}
#include <cuda_runtime.h>
#include <cstddef>
#include <new>
#include <limits>

template <typename T>
struct CudaAllocator {
    using value_type = T;
    using size_type = std::size_t;
    using difference_type = std::ptrdiff_t;

    CudaAllocator() noexcept = default;

    template <typename U>
    CudaAllocator(const CudaAllocator<U>&) noexcept {}

    T* allocate(std::size_t n) {
        if (n > std::numeric_limits<std::size_t>::max()
                / sizeof(T))
            throw std::bad_alloc();
        T* ptr = nullptr;
        auto err = cudaMalloc(&ptr, n * sizeof(T));
        if (err != cudaSuccess || !ptr)
            throw std::bad_alloc();
        return ptr;
    }

    void deallocate(T* p, std::size_t) noexcept {
        if (p) cudaFree(p);
    }

    template <typename U>
    bool operator==(const CudaAllocator<U>&) const {
        return true;
    }
    template <typename U>
    bool operator!=(const CudaAllocator<U>&) const {
        return false;
    }
};
\end{lstlisting}

\subsection{CudaVector 封装}

与方式一不同，\texttt{CudaVector} 现在是一个普通类，
对象本身由标准 \texttt{new} 分配到主机堆上：

\begin{lstlisting}
#include <vector>
#include <algorithm>

template <typename T>
class CudaVector {
public:
    CudaVector() = default;

    explicit CudaVector(std::size_t n)
        : device_data_(n) {}

    CudaVector(std::size_t n, const T& val)
        : device_data_(n, val) {}

    // 无需自定义 operator new/delete
    // 对象本身在主机堆上分配（默认行为）

    void push_back(const T& val) {
        // 通过 staging 上传
        T host_val = val;
        device_data_.push_back(host_val);
    }

    T read(std::size_t i) const {
        // device_data_ 的元素在显存中
        // 需要通过 cudaMemcpy 读取
        T val;
        cudaMemcpy(&val,
                   device_data_.data() + i,
                   sizeof(T),
                   cudaMemcpyDeviceToHost);
        return val;
    }

    void write(std::size_t i, const T& val) {
        cudaMemcpy(device_data_.data() + i,
                   &val, sizeof(T),
                   cudaMemcpyHostToDevice);
    }

    std::size_t size() const {
        return device_data_.size();
    }

    // 供 kernel 使用
    T* device_data() { return device_data_.data(); }

private:
    std::vector<T, CudaAllocator<T>> device_data_;
};
\end{lstlisting}

\subsection{使用方式}

\begin{lstlisting}
int main() {
    // 对象在主机堆上（普通 new）
    auto* vec = new CudaVector<float>(1024, 1.0f);

    // kernel 调用方式相同
    int threads = 256;
    int blocks = (1024 + threads - 1) / threads;
    scale_kernel<<<blocks, threads>>>(
        vec->device_data(),
        static_cast<int>(vec->size()));
    cudaDeviceSynchronize();

    float val = vec->read(0);
    delete vec;   // 普通 delete，触发析构
    return 0;     // 析构中 CudaAllocator 释放显存
}
\end{lstlisting}

\section{两种方式的关键差异}

\begin{compare}[title=new/delete 重载 vs Allocator 封装]
\textbf{方式一：类级别 operator new/delete}
\begin{itemize}
    \item \texttt{CudaVector} 对象本身在 CUDA 统一内存中。
    \item 对象控制块（\texttt{device\_ptr\_}、\texttt{size\_} 等）可被 GPU 直接访问。
    \item 需要显式管理显存数据缓冲区的生命周期（构造函数/\allowbreak 析构函数中 \texttt{cudaMalloc}/\texttt{cudaFree}）。
    \item 适合需要在 GPU kernel 中访问向量元数据（如 size）的场景。
    \item \texttt{delete} 表达式同时释放统一内存中的对象和显存中的数据。
\end{itemize}

\textbf{方式二：自定义 CudaAllocator}
\begin{itemize}
    \item \texttt{CudaVector} 对象本身在主机堆上（标准 \texttt{new}）。
    \item 仅元素数据通过 \texttt{CudaAllocator} 存储在显存中。
    \item 利用 \texttt{std::vector} 管理显存缓冲区的生命周期（RAII），代码更简洁。
    \item GPU kernel 无法直接访问向量的 \texttt{size()}（需传参）。
    \item 对象析构时，\texttt{vector} 析构自动调用 \texttt{CudaAllocator::deallocate} 释放显存。
\end{itemize}
\end{compare}

\section{选择建议}

\begin{bestpractice}[title=CUDA 封装方式选择]
\textbf{优先使用 Allocator 方式}，除非有明确的理由需要重载 operator new。
具体考量：

\begin{itemize}
    \item \textbf{代码量}：Allocator 方式复用 \texttt{std::vector} 的内存管理逻辑，
          自定义代码量少约 50\%。operator new 方式需要手写全部缓冲区管理。
    \item \textbf{RAII 安全性}：Allocator 方式中，\texttt{vector} 析构自动释放显存。
          operator new 方式需要确保所有退出路径（包括异常）都正确释放。
    \item \textbf{GPU 元数据访问}：如果 kernel 需要读取 \texttt{size()} 等元数据，
          方式一（统一内存）可以零拷贝访问；方式二需要将 size 作为 kernel 参数传入。
    \item \textbf{性能控制}：方式一使用 \texttt{cudaMallocManaged} 有 page fault 开销。
          对延迟敏感的代码，应在方式二中将元数据和数据分别放在主机和显存。
    \item \textbf{与 STL 算法的兼容}：方式二的内部 \texttt{vector} 可以与
          \texttt{thrust} 等 GPU STL 库直接互操作。方式一需要额外的适配。
\end{itemize}
\end{bestpractice}

\section{混合方案}

在高性能场景中，可以组合两种机制：用 Allocator 管理显存数据，
用 placement new 将对象控制块放到固定地址的对齐内存中：

\begin{lstlisting}
#include <new>

// 对齐的主机内存（适合 pinned memory 传输）
alignas(64) static char control_block[256];
static std::size_t cb_offset = 0;

// 从预分配的对齐内存中分配
template <typename T, typename... Args>
T* create_pinned(Args&&... args) {
    std::size_t aligned =
        (cb_offset + alignof(T) - 1)
        & ~(alignof(T) - 1);
    void* ptr = &control_block[aligned];
    cb_offset = aligned + sizeof(T);
    return new (ptr) T(
        std::forward<Args>(args)...);
}

// 使用:
auto* vec = create_pinned<CudaVector<float>>(1024, 1.0f);
// vec 在 64 字节对齐的主机内存中
// vec->device_data() 指向显存
// 可以用 cudaMemcpyAsync + pinned memory 实现
// 零拷贝异步传输
\end{lstlisting}

\begin{warning}[title=混合方案的析构]
使用 placement new 创建的对象不能用 \texttt{delete} 释放。
必须手动调用析构函数，然后自行管理底层内存的回收：
\begin{lstlisting}
vec->~CudaVector();   // 手动析构，释放显存
// control_block 内存由静态数组管理，无需释放
\end{lstlisting}
\end{warning}

\end{document}
